\chapter{Critical Limitations and Hardware Realities}
\label{chap:limitations}

\section{Scientific Integrity: No Quantum Advantage Claim}
A foundational commitment of this research report is absolute scientific integrity. We state plainly: \textbf{this work does not demonstrate, nor does it claim, quantum advantage or quantum speedup.}

For the evaluated $N = 5$ instance, classical exact brute-force search enumerates all $2^5 = 32$ states in $0.49$ milliseconds on a standard desktop CPU, guaranteeing identification of the globally optimal solution. In contrast, noiseless statevector simulation of QAOA on the same processor requires $1.42$ seconds at depth $p=1$ and $9.77$ seconds at depth $p=3$---nearly four orders of magnitude slower. Classical industrial solvers such as Gurobi or CPLEX routinely solve mixed-integer quadratic programming (MIQP) instances with thousands of binary variables in seconds using branch-and-bound and cutting-plane algorithms.

The value of this implementation lies in establishing a rigorous, reproducible, mathematically verified reference pipeline for constrained combinatorial optimization in Qiskit, demonstrating how quantum algorithms handle constraints, and evaluating the physical bottlenecks that must be overcome before quantum utility is achievable.

\section{Physical NISQ Hardware Bottlenecks}
Executing the Quantum Resource Optimizer on physical gate-model superconducting processors introduces severe physical noise sources absent in idealized statevector simulation \cite{preskill2018quantum, sachdeva2024quantum}:

\subsection{Two-Qubit Gate Error Accumulation}
On contemporary superconducting transmon processors (such as IBM Quantum Eagle and Heron chips), single-qubit gate fidelities exceed $99.9\%$, but two-qubit entangling CNOT/ECR gates suffer from error rates of $\epsilon_{\text{CX}} \approx 0.5\% - 1.0\%$.

For an unconstrained graph with $N$ vertices, the number of two-qubit gates per QAOA layer is $\binom{N}{2} = \frac{N(N-1)}{2}$. For $N = 5$, each layer requires 10 entangling operations, meaning a $p=3$ circuit executes 30 two-qubit gates. In a first-order depolarizing noise model, circuit fidelity degrades exponentially:
\begin{equation}
\mathcal{F} \approx (1 - \epsilon_{\text{CX}})^{N_{\text{CX}}} \approx (1 - 0.008)^{30} \approx 78.5\%
\end{equation}
As $N$ scales to realistic portfolio sizes ($N = 50$), each layer requires $\binom{50}{2} = 1225$ two-qubit gates. At $p=3$, over 3600 CNOTs would be required, collapsing circuit fidelity to zero ($\mathcal{F} \ll 10^{-12}$) and rendering raw physical measurements indistinguishable from uniform white noise.

\subsection{Decoherence \texorpdfstring{($T_1, T_2$)}{(T1, T2)} Constraints}
Superconducting transmons possess finite energy relaxation times ($T_1 \sim 150 - 300\,\mu\text{s}$) and phase coherence times ($T_2 \sim 100 - 200\,\mu\text{s}$). A typical two-qubit cross-resonance gate requires $\sim 300 - 500\,\text{ns}$. A circuit with deep layers and sequential entangling operations quickly exhausts the coherence budget, causing the quantum register to decay toward a maximally mixed thermal state before measurement.

\section{Lattice Connectivity and SWAP Network Overhead}
A critical structural challenge in quantum portfolio optimization is graph density. Because every pair of assets shares a covariance value $\Sigma_{ij}$ and a quadratic budget penalty $2\lambda$, the Ising Cost Hamiltonian graph is a \textit{complete graph} ($K_N$):
\begin{equation}
\text{Degree}(v_i) = N - 1, \quad \forall i
\end{equation}

However, physical superconducting QPUs are constrained to planar 2D lattices with low qubit connectivity. IBM Quantum processors employ a \textit{heavy-hex} lattice architecture with node degrees strictly limited to 2 and 3.

To execute two-qubit interactions between physically non-adjacent qubits on a heavy-hex chip, the quantum compiler must insert SWAP networks. Each logical two-qubit gate between distant qubits requires routing via intermediate SWAP gates, where each SWAP decomposes into 3 physical CNOT gates. Compiling a dense $K_N$ graph onto a heavy-hex topology multiplies the physical CNOT depth by an additional factor of $3\times$ to $5\times$, drastically accelerating decoherence and gate error accumulation.

\section{Variational Optimization Landscapes: Barren Plateaus}
The classical optimization of variational angles $(\vec{\gamma}, \vec{\beta})$ faces theoretical hurdles as instance sizes expand. McClean et al. \cite{mcclean2018barren} proved that in generic parameterized quantum circuits, the variance of the gradient of the expectation value vanishes exponentially in the number of qubits:
\begin{equation}
\text{Var}_{\vec{\theta}}\left[ \frac{\partial \langle \hat{H}_C \rangle}{\partial \theta_k} \right] \in \mathcal{O}\left( \frac{1}{2^N} \right)
\end{equation}

For large $N$, the optimization landscape becomes an exponentially flat ``barren plateau'' punctuated by infinitesimally narrow global minima basins. In the presence of hardware noise, this phenomenon is further exacerbated: noise-induced barren plateaus flatten the landscape regardless of circuit depth, trapping gradient-based and simplex-based classical optimizers like COBYLA in random local configurations.
